\section{The critical horizontal rate and the return to wide intervals}
\label{sec:critical-rates}

We complete the finite-narrowness argument in the horizontal model of
\Cref{prop:horizontal-model}.  All depths in this section are relative to
the selected true parent \(Q_r\).  Thus
\[
 a_s=N^{-s},\qquad h_s=N^{-ks},\qquad 0<k=\ell\le1,
 \qquad m_{Q_s}=H_s\nu^\flat(Q_s)
       \sim_{\log,N}h_s^4a_s^d,
\]
where \(0\le d<1\), \(H_s\) is the true block length in the normalized
parent, and \(\nu^\flat\) is its trajectory prior.  The constants denoted
by \(K\) are subpower in these tangent units.  We retain the full true
assignment weights and the earlier witnesses when restricting the current
incidence law.  Restoring a counted mass to the initial tangent parent
multiplies it by \(q_r\nu(Q_r)\), up to the stated subpower factors.

Write \(z=(t,Z)\), \(V_1=(1,v)\), and \(J(t,Z)=(-Z,t)\).  The
normalized scalar \(x=(w-L)/a_r\) is affine on each participating
trajectory, with \(|x|+|x'|\le K\).  The derivative \(x'\) in this section
always uses the full normalized \(Q_r\)-time variable.  We use the
horizontal equation, affine slope estimates, and coefficient counts of
\Cref{a06:horizontal-equation,a06:affine-slope-bounds,a06:horizontal-coefficient-count}.  In particular, with
\(P=(p,q_0)=(Z(0),Y(0))\),
\begin{equation}
 Z=p+tv,\qquad Y=q_0-tp+O(N^{-e_0})
 \label{a07:parameter-reconstruction}
\end{equation}
holds throughout the participating lines for a fixed \(e_0>0\).
The pure and joint caps, together with the end-history palette bounds of
\Cref{prop:tangent-palette,a03:finite-state-conditioning}, are upper bounds on the unnormalized
restricted measures.  The palette \(\pi\) of \(v\) has its interval
nonconcentration bound.  These conventions will be used even after an
individual descendant or predictor is selected.

At the critical rate \(2k=1\), we fit the derivative of \(x\)
and then count its remaining trajectory-dependent intercepts. Three
successive changes of trajectory show that intercept bins carrying
most of the mass each occupy a substantial part of the three-dimensional base; the true descendant
counts then force one heavy bin. For \(2k\ne1\), we instead construct
wide projected intervals and apply \Cref{prop:wide-interval}. Below
the critical rate this requires full-time predictors on parameter
boxes of width \(h_s^2\); \Cref{prop:bootstrap} reaches that scale
from the coefficient counts of the preceding section.

Here is the elementary deletion principle used to make several finite
families of labels available at once.

\begin{lemma}[Simultaneous marginal floors]
\label{a07:simultaneous-floors}
Let a finite measure of mass \(M\ge K^{-1}\) carry finitely many finite
label families.  Suppose that each family partitions its incidences, up
to a subpower multiplicity, and that its proposed positive floor
numerators \(T_j\) have total at most \(K\), summed over all families.
For a sufficiently large subpower \(K_1\), one may retain at least
\(9M/10\) so that every remaining label has current mass at least
\(T_j/K_1\).  Earlier label witnesses and trajectory priors may be kept
without restriction as reference objects.
\end{lemma}

\begin{proof}
Use the fixed-threshold deletion argument of
\Cref{a01:witness-budgets}, with thresholds \(T_j/K_1\) and
\(K_1\) chosen so that \(\sum_jT_j/K_1<M/10\), including the
subpower multiplicities.  Each deleted label becomes permanently empty,
so the process terminates after finitely many deletions and charges each
label at most once; the total charge is less than \(M/10\).
Only the current measure is restricted; earlier witnesses and trajectory
priors remain unchanged.  For increasing finite grids, choose their sizes
slowly enough that the multiplicities and \(K_1\) remain subpower.
\end{proof}

\subsection{The critical rate}

\begin{proposition}[Critical-rate improvement]
\label{prop:critical-rate}
If \(2k=1\), every substantial residual of the prepared horizontal
problem yields a time-improvement candidate in the sense of
\Cref{prop:candidate-upgrade}.  More precisely, for a sufficiently small
fixed \(s>0\) and any fixed \(0<c<ks/2\), one can choose a true
\(Q_s\) and an affine base predictor whose hidden-variable test has
width \(K a_r a_{s+c}\), hidden derivative one, and counted mass at
least
\[
 K^{-1}q_r\nu(Q_r)m_{Q_s}N^{-dc}
\]
in the initial tangent parent.  Its fit holds throughout the selected
block on the assigned trajectories.
\end{proposition}

\begin{proof}
\par\smallskip\noindent\textbf{An affine primitive for the trajectory derivatives.}\quad

We first find an affine base function \(F_*\) whose derivative
follows \(x'\) on a dense family in one \(Q_s\). Integrating
leaves an intercept for each trajectory. At the finer width
\(a'=a_sN^{-c}\), choose a sufficiently fine base mesh of side
\(\omega\). True descendant counts bound the number of occupied
pairs of base cells and intercept bins by
\(K\omega^{-3}(a_s/a')^d\). We will show that bins carrying
almost all the fitted mass each occupy at least
\(K^{-1}\omega^{-3}\) base cells. Their number is therefore at
most \(K(a_s/a')^d\), forcing the desired heavy intercept.


Fix \(s>0\) sufficiently small that all required multiples of \(s\) lie in the horizontal preparation range and \(100s<e_0\), put \(a=a_s=h_s^2\), and apply \Cref{a06:horizontal-coefficient-count} with \(D=a\). Direction and time meshes \(d_0,\eta\) can be e.g. of sizes \(N^{-100s}\); use finer fixed-power meshes for summations/conditionings as needed. Fix also \(0<c<ks/2\); include depths \(s+c,3s\) in the preparations before zooming in a \(Q_s\).

At an occurrence in \(Q_s\) write \(H=H_s,\ \widehat z=(z-z_c)/H\) now with a box center as in \Cref{prop:horizontal-model}. Given \(Q_s,\widehat z_\xi,v_\xi\) for sufficiently fine fixed power \(\xi\ll a\) (finer also than needed for \(d_0,\eta\)), there are only \(K\) options for the \(x'\)-bin of width \(a\), with prime still in full \(Q_r\)-time units. Indeed \(p\) is recovered using \(Z=p+t v\); \(Y\) is bounded to \(K H^2\) accuracy using \(Q_s,z\); hence \(q_0\) is predicted using \(q_0=Y+t p+O(N^{-e_0})\), with only \(K\) options for \(P_D\) altogether, since \(D=a=h_s^2\). Likewise the \(L_s\)-lists restrict \(x_{a}\) with only \(K\) options. The time subbin and \(d_0\) are known up to neighbors. Thus the
group, time subbin and possible \(x_a\)-bins are available, and
\Cref{a06:horizontal-coefficient-count} gives the derivative list.

We next turn the derivative lists into a list of possible rows.
On this law,
\[
 x'=G_*\cdot V_1+O(K a),\qquad G_*=\nabla_z L_{3s}+(\partial_Y L_{3s})Jz,\qquad |G_*|\le K ,
\]
by \Cref{a06:horizontal-equation,a06:affine-slope-bounds}, and differentiating the affine fits along lines (error \(\le K a_{3s}/h_{3s}\) by the fitted-block bounds).

Use typical states of \((Q_s,\widehat z_\xi,[G_*]_a)\), with conditional velocity-bin cap \(K\pi\) after light-entry and light-conditioning-state deletion. Indeed this appended state costs only \(K\) options given sufficiently deep histories: \(L_{3s}\) has only \(K\) entries for its label, and \Cref{a06:affine-slope-bounds} here controls row variation. The end-label domination with earlier masses summing inside \(Q_r\) to \(\le K q_r\nu_{Q_r}\) thus applies by \Cref{a03:finite-state-conditioning}.

Each surviving such state has supporting pairs of velocity bins, separated by at least \(1/K\), of palette product weight \(\ge1/K\). Use conditional domination and the angular interval bound, choosing the inverse-subpower separation threshold sufficiently small relative to prior subpower losses. For any fixed such pair over \(Q_s,\widehat z_\xi\), the two \(x'\) evaluation lists constrain the compatible row-bin centers to \(K\) options by linear interpolation (evaluation errors \(K a\) using bounded rows and velocities). Summing this multiplicity bound with the palette product weights gives
only \(K\) occurring \(a\)-bins for \(G_*\) per
\(Q_s,\widehat z_\xi\).

The remaining input for \Cref{a06:row-fit} is the joint cell cap.
Pick one \(Q_s\) with dense success relative to \(m_{Q_s}\).
Integrate the original joint caps on much finer base meshes, at
\(v_\xi\), over the transverse range \(K H^2\) and the \(K\)
affine predictions for \(x\) at accuracy \(K a\). Before division
by \(m_{Q_s}\), the bound is \(K(H\xi)^2 H^2 a^d\pi(v_\xi)\);
afterwards it is \(K\xi^2\pi(v_\xi)\).

In centered coordinates horizontal velocity remains \(V_1\) and time has length one (constant translations are harmless). Apply \Cref{a06:row-fit} to \(T=G_*,z=\widehat z,E=a\), with line constant \(x'\) and the block-conditioned prior. On dense mass relative to \(m_{Q_s}\), \(G_*=B+\lambda J\widehat z+O(K a)\), \(|B|+|\lambda|\le K\). To realize this fitted row along the trajectories, observe that the affine function
\[
 F_*(z,Y)=(B-\lambda Jz_c/H)\cdot z+\lambda Y/H
\]
has slopes at most \(K/H\) and satisfies
\[
 |x'-\tfrac{d}{dt}F_*(z,Y)|\le K a.
\]
Indeed use \Cref{a06:horizontal-equation} with \(N^{-e_0}/H\ll a\); these last two trajectory derivatives are constant.

\par\smallskip\noindent\textbf{Intercepts and the upper cell count.}\quad
Take \(a'=aN^{-c}\).  For each fitted indexed trajectory \(l\),
let \(c_l=(x-F_*)(t_{\rm mid})\) be its exact midpoint value,
and let \(C=[c_l]_{a'}\) denote its bin.  The exact value obeys
\[
 x(t)-F_*(z(t),Y(t))=c_l+O(Ka h_s),
 \qquad Ka h_s\ll a',
\]
throughout the block, since \(c<ks/2\).  Thus a sufficiently precise
current base position and \(x\)-value determine only boundedly many
possible \(C\)'s.  Replacing \(c_l\) by its bin center incurs
\(O(a')\), which is included in the eventual fit error.

In centered coordinates \((\widehat z,\widehat Y)\), with plane \(J z_c\) subtracted and transverse scaling by \(H^2\), use a base mesh \(\omega\ll a'\), say \(N^{-5s}\). The total number of (base cell, intercept bin) pairs active here costs
\[
 K\omega^{-3}(a/a')^d .
\]
Indeed the number of true descendants \(Q_{s+c}\), over all their subblocks,
is at most
\[
 K\left(\frac{a}{a'}\right)^d
   \left(\frac{h_s}{h_{s+c}}\right)^4,
\]
by the time-length ratios and disjoint nested trajectory assignments.
Each descendant's domain has volume at most
\(K(h_{s+c}/h_s)^4\) in centered \(Q_s\) coordinates,
including mesh enlargement. To see this, use its horizontal scale
and transverse square scale at a meeting center. Its centered plane
has tilt at most \(K\) in \(Q_s\) units, and \(\omega\)
is finer by a fixed power than its relative widths.
Multiplying this volume bound by \(O(\omega^{-3})\) pays for the
cells in each box. Summing over boxes cancels the two fourth-power
scale ratios. By \Cref{a06:affine-slope-bounds} and the slopes of \(F_*\), its \(L_{s+c}\)-entries specify only \(K\) intercept bins per such cell.

\par\smallskip\noindent\textbf{Three horizontal traversals and the lower cell count.}\quad
We now prove the complementary estimate: intercept bins carrying
almost all the fitted mass each occupy at least
\(K^{-1}\omega^{-3}\) base cells.  Together with the preceding
upper count, this will bound the number of such bins by
\(K(a/a')^d\).

The switching state records \(C\), centered base position, and
\(x\), the latter two on meshes much finer than \(\omega\).
It has only \(K\) possibilities per sufficiently deep true
history.  Indeed \Cref{lem:horizontal-affine} gives this conclusion
for the fine base and residual bins, even beyond the accuracy of the
horizontal equation.  Those bins determine only \(K\) possible
\(C\)'s, because \(F_*\) has the stated slope bounds and
\(x-F_*=c_l+O(Ka h_s)\), with \(Ka h_s\ll a'\).

Apply the reference-weight pruning of \Cref{a06:summable-pruning}
to these appended states inside the selected \(Q_s\).  The
earlier end-history weights satisfy
\(\sum_eR_e\le Km_{Q_s}\), and the velocity numerator for
history \(e\) is bounded by \(K\pi(b)R_e\).  Deleting
light history--state entries, and then any light aggregate states
needed for the conditioning, retains a law \(\mu_{\rm last}\)
with
\[
 m_{\rm last}:=\mu_{\rm last}(\mathrm{all})\ge m_{Q_s}/K
\]
and conditional velocity-bin weights at most \(K\pi(b)\).
All reference weights and earlier full-domain witnesses are unchanged.
Set \(\widetilde\mu=\mu_{\rm last}/m_{\rm last}\).
This probability is bounded by \(K\) times the true trajectory
prior in \(Q_s\) times uniform centered time.

Call a trajectory good when its \(\widetilde\mu\)-marginal
density relative to that prior is at least \(1/K_1\), where
\(K_1\) is a sufficiently large subpower.  Bad trajectories
carry negligible mass, and the conditional time density on every
good trajectory is at most \(K\), after enlarging \(K\).

Starting with an event of law \(\widetilde\mu\), perform the
following two draws three times.  First draw a fresh event conditional
on the current switching state.  Then, conditional on its entire
indexed trajectory, draw a new event on that trajectory.  Both
conditional resampling operations preserve \(\widetilde\mu\).
They also preserve \(C\), which is both trajectory data and part
of the switching state.  Denote the rounded velocity on the \(i\)th
fresh trajectory by \(c_i\), and its new centered time by
\(\tau_i\).

Retain for the estimate only paths whose three fresh trajectories are
good and for which
\[
 |c_2-c_1|,\ |c_3-c_2|,\ |\tau_2-\tau_1|\ge1/K'.
\]
Here \(K'\) can be chosen as a sufficiently large subpower so
that the total failure probability is \(o(1)\).  This follows
from stationarity, the conditional velocity bound and the palette's
interval bound, and the time-density bound on good trajectories.
We restrict paths without renormalizing any transition kernel.
In particular, conditional on the full past, the subkernel of
\((c_i,\tau_i)\) on good fresh trajectories is bounded by
\(K\pi\otimes d\tau_i\).  Successive use of this bound
therefore dominates the joint control--time law by
\(K\pi^{\otimes3}\otimes d\tau_1d\tau_2d\tau_3\).
The bound holds before imposing endpoint or later regularity tests.

Up to error \(o(\omega)\), the three traversals follow
\[
 \frac{d\widehat Z}{d\widehat t}=c_i,
 \qquad
 \frac{d\widehat Y}{d\widehat t}=-\widehat Z+\widehat t c_i.
\]
The refresh changes the start of each traversal only within the tiny
base mesh.  The horizontal equation contributes error
\(O(KN^{-e_0}/H)\), and velocity rounding contributes still
less.  Take the meshes sufficiently fine and \(s\) sufficiently
small that the accumulated error over three traversals is
\(o(\omega)\).  Signed time increments are allowed.

Fix the starting event \((\tau_0,Z_0,Y_0)\) and the control
triple, and put \(D_j=c_{j+1}-c_j\).  Exact integration gives
\begin{align}
 Z_f&=Z_0+c_1(\tau_3-\tau_0)
        +\sum_{j=1}^2D_j(\tau_3-\tau_j),\notag\\
 Y_f&=Y_0+(-Z_0+\tau_0c_1)(\tau_3-\tau_0)
        +\sum_{j=1}^2D_j\tau_j(\tau_3-\tau_j).
 \label{a07:critical-endpoint}
\end{align}
A switch at \(\tau_j\) changes the transverse velocity by
\(\tau_jD_j\).  Consequently
\begin{align}
 \det\frac{\partial(\tau_3,Z_f,Y_f)}
              {\partial(\tau_1,\tau_2,\tau_3)}
 &=\det\begin{pmatrix}
 -D_1&-D_2\\
 D_1(\tau_3-2\tau_1)&D_2(\tau_3-2\tau_2)
 \end{pmatrix}\notag\\
 &=2D_1D_2(\tau_2-\tau_1).
 \label{a07:critical-jacobian}
\end{align}
This determinant is bounded below by an inverse subpower on the
retained paths.  At fixed \(\tau_3,Z_f\), the internal times
satisfy one nontrivial linear equation.  On that line the equation
for \(Y_f\) has degree at most two and cannot be constant at a
regular solution.  Thus there are at most two regular preimages.
The change-of-variables formula bounds the time-triple measure
mapping into an enlarged base cell by \(K\omega^3\), uniformly
in the fixed starting event and control triple.  Integrating against
\(\pi^{\otimes3}\) and using the preceding subkernel bound gives
\begin{equation}
 \Pp(\text{regular endpoint in a specified base cell}
             \mid\text{start})\le K\omega^3.
 \label{a07:critical-cell-cap}
\end{equation}
The simulation error requires only a bounded enlargement of the cell.

Write \(q_C=\widetilde\mu(C)\), and let \(e_C\) be the
failure probability averaged over the starting law conditional on
\(C\).  Since the total failure probability is
\(\sum_Cq_Ce_C=\eta=o(1)\), the bins
\(\mathcal B=\{C:e_C\le1/2\}\) have total mass at least
\(1-2\eta\).  For each such bin, integrate
\Cref{a07:critical-cell-cap} over its conditional starting law.
Every endpoint still has intercept bin \(C\), while regular
paths have probability at least \(1/2\).  Hence
\[
 \#\{\text{base cells occupied by }C\}
       \ge\frac1{2K\omega^3},\qquad C\in\mathcal B.
\]
This is the lower cell count on bins carrying almost all the fitted
mass.


\par\smallskip\noindent\textbf{The heavy candidate in the original mass units.}\quad  The upper count of \((\text{base cell},C)\) pairs and the lower count
for \(C\in\mathcal B\) give \(\#\mathcal B\le K(a/a')^d\).
Since \(\sum_{C\in\mathcal B}q_C\ge1-2\eta\), one of them has
\[
 \mu_{\rm last}(C)=m_{\rm last}q_C
 \ge K^{-1}m_{Q_s}(a'/a)^d
 =K^{-1}m_{Q_s}N^{-dc}.
\]
This is mass on the original residual in horizontal-problem units,
before probability normalization. Its \(x-F_*\) offset test matches throughout the fitted
block at width \(K a'\). Restoring \(w\) gives the affine test
\(w-L-a_r(F_*+\text{offset})\) in the initial tangent parent,
with hidden derivative one and width \(K a_r a'\).
Restoring mass multiplies, up to \(K\), by \(q_r\nu_{Q_r}\),
yielding the required counted mass
\(K^{-1}q_r\nu_{Q_r}m_{Q_s}N^{-dc}\).
This is the time-improvement candidate at initial depth \(r+s\).

\end{proof}

\subsection{Full-time parameter estimates below the critical rate}

The coefficient count gives predictors at thickness \(a_s\) on
parameter boxes of width \(a_s\). For the return to wide intervals
we need the larger width \(h_s^2=a_s^{2k}\). We record the
intermediate parameter widths as follows.

\begin{definition}[Full-time parameter estimates]
\label{a07:bootstrap-property}
Suppose \(2k<1\). Choose \(s_{\max}>0\) sufficiently small
compared with \(e_0\), allowing true-chart preparations up to the
required fixed multiples of \(s_{\max}\). Fix a common velocity
grid \(d_0\) of width \(\chi\le N^{-e_0}\), a fixed power
up to dyadic rounding.

For \(2k\le u\le1\), property \(\mathcal G(u)\) means the
following. For every fixed \(0<s<s_{\max}\) and every dense
input, a further dense restriction and subsequence have full-time
labels in the groups \((d_0,P_D)\), where \(D=a_s^u\).
Their trajectory assignments are disjoint, their individual active
masses are at least
\[
 K^{-1}\pi(d_0)D^2a_s^d,
\]
and each label has a predictor
\[
 x=H(t,P)+O(Ka_s)
\]
throughout unit time on its assigned trajectories. The function
\(H\) is affine in \(P\), with coefficients affine in time,
and is bounded by \(K\) on the group's scaled parameter box
over unit time. The counted incidences may be restricted in time.
All masses and depths use the fixed horizontal-model conventions;
dyadic rounding and subpower factors are allowed as before.
\end{definition}

\begin{lemma}[Comparison with full-time predictors]
\label{a07:full-time-comparison}
Suppose \(2k<1\), \(2k\le u\le1\), and
\(\mathcal G(u)\) holds. Fix a sufficiently small depth \(s_0\), and put
\(a_0=a_{s_0}\), \(D_0=a_0^u\), and \(G=(d_0,P_{D_0})\).
First prepare true labels and affine entries at \(s_0\) and at
any prescribed finite set of nearby deeper depths. Then apply
\(\mathcal G(u)\) to the resulting dense input, obtaining
full-time labels \(H_0\) at \(s_0\). On a further dense
restriction, the following conclusions hold.

For each \(Q_{s_0}\) time block \(I\), write
\[
 {\cal L}_{s_0}(t,P)
   =L_{s_0}(t,p+t v_{d_0},q_0-tp).
\]
Every retained pairing of \(H_0\) with \((Q_{s_0},L_{s_0})\)
satisfies
\[
 \|{\cal L}_{s_0}-H_0\|\le Ka_0
\]
on the whole product of \(I\) and its parameter box. Given
\(H_0,I\), at most \(K\) such true-label entries occur.
For each prescribed deeper depth \(s\), the paired true affine
entries also have stored full-domain witnesses of mass at least
\(m_{Q_s}/K\), and
\(\|L_s-L_{s_0}\|\le Ka_0\) on the whole child box.
The witnesses and full trajectory assignment priors are kept when
the current incidence law is further restricted.
\end{lemma}

\begin{proof}
For nested true labels we use \Cref{a07:simultaneous-floors} to retain full-domain witnesses for paired affine entries at \(s_0\) and any nearby required deeper relative depths \(s\). Prepare as in \Cref{a06:affine-slope-bounds} before the small slices or full-time selections: discard light pairs compared with \(m_{Q_s}=H_s\nu^\flat(Q_s)\), using the \(\le K\) choices of matched subentries per true child label including its ancestor. Each remaining pair has earlier witnesses together near both fits, weighing \(\ge m_{Q_s}/K\). Thus \(L_s-L_{s_0}\) costs \(K a_0\) in norm throughout the child box by affine Remez, using inverse-subpower base coverage from the child mass and density cap there. We can prepare all such data first on sufficiently fine fixed finite grids, increasing sufficiently slowly.

Slice by \(G=(d_0,P_{D_0})\) and \(Q_{s_0}\)'s time block \(I\);
put \(W_G=\pi(d_0)D_0^2\). Both lists \(H_0\) and
\((Q_{s_0},L_{s_0})\) have typical individual marginal floor
\(T_{GI}/K,\ T_{GI}=|I|W_G a_0^d\) there. Indeed the thresholds
summed for \(H_0\) use the earlier full-time label floors and
\(\sum_I|I|=1\).

For each \(Q_{s_0},d_0\) there are at most
\(K h_{s_0}^3/D_0^2\) parameter cells of \(P\) to test. Use true
assigned positions of contributing trajectories at the time midpoint,
within horizontal scale \(K h_{s_0}\) and tilted transverse scale
\(K h_{s_0}^2\). The horizontal equation error on these lines is
constant in velocity through the block by
\Cref{a06:horizontal-equation}. Thus, up to
\(K(\chi+N^{-e_0})\ll D_0\), the spatial positions come from
\(P\mapsto(p+t v_{d_0},q_0-tp)\), where \(v_{d_0}\) is the binned
velocity; this has determinant one and distortion bounded by \(K\).
A meeting parameter bin fits the \(K\)-enlarged range since
\(D_0\le K h_{s_0}^2\).

Sum the thresholds using this count, the \(K\) affine subentries,
\(\sum\pi(d_0)=1\) and \(\sum m_{Q_{s_0}}\le1\). Thus for each
list the sum of thresholds \(T_{GI}\) costs \(\le K\), and light
individual entries can be deleted (keeping their pre-deletion
counterparts as witnesses for the heavy ones). The assignments counted
within each list on these domains are disjoint up to the
\(K\)-subentry losses.

On this product domain, \(H_0\) and \({\cal L}_{s_0}\) both
match \(x\) to \(Ka_0\) at paired points, by
\Cref{a06:horizontal-equation,a06:affine-slope-bounds}. We use
\Cref{a02:local-graph-matching} to upgrade this agreement to the
whole domain. The same calligraphic notation will be used for
true affine entries at other depths.

For either index on sufficiently fine \((t,P)\) meshes, index-base marginals obey the upper bound \(K T_{GI}\) times box-uniform, by the joint cap at \(d_0,x_{a_0}\) near its graph. Indeed use the map \(P\mapsto(p+tv_{d_0},q_0-tp)\) to true spatial positions up to \(K(\chi+N^{-e_0})\), summing over much finer true base cells. Partition the time and parameter domains internally with absolute side steps on the order of \(N^{-C s_{\max}}\), with sufficiently large fixed \(C\) and still \(C s_{\max}\) well below \(e_0\). These dominate the reconstruction error, and errors inside such cells in both graph evaluations are negligible compared with \(a_0\). The true \(L_{s_0}\) slopes are bounded by \(K\) here; \(H_0\) has polynomial norm \(\le K\) on the full-time \(D_0\)-scaled group. Thus jitter the paired base point together within its cell to realize the density ceilings.

If a power-excess norm discrepancy occurs on dense mass, pigeonhole a comparable discrepancy \(M_0\), larger than \(a_0\) by a power (the occurring graph norms themselves cost \(K\) by those bounds and an active value). On each domain group the two lists into combinations of neighboring polynomial jet bins at size \(M_0\) in domain-scaled variables; paired discrepancies of that size permit only boundedly many neighbors (up to \(K\)). Within such groups both list counts are \(\le K(M_0/a_0)^d\): sum their individual earlier \(T_{GI}/K\) floors under the cap near a common graph representative at \(K M_0\)-thickness, integrating the joint cap over that domain. The weights \(T_{GI}\) summed per index including repetitions in neighboring groups still cost \(\le K\) over all domains. Choose a group of dense bad-pair mass relative to the summed index weights. Subtract a common representative there. Jittering preserves pairing as stated, and after normalizing the law each index-base marginal fits relative to the normalized index weights up to \(K\). \Cref{a02:local-graph-matching} with \(R=M_0/a_0\), applied to these polynomial graphs on the unit-scaled box, gives impossibility. Thus send the power tolerance slowly to zero and prune failures. At the remaining \(K a_0\) norm discrepancy, for fixed \(H_0,I\) the individual earlier floors for the possible other indices all count near its graph at \(K a_0\)-thickness by norm closeness, giving the asserted \(K\)-list by the cap again.

\end{proof}

\begin{proposition}[Parameter bootstrap below the critical rate]
\label{prop:bootstrap}
If \(2k<1\), there is a fixed \(s_{\max}>0\) for which
\(\mathcal G(2k)\) holds.  It can be applied anew to any dense input
sequence and any subsequence.  Any prescribed finite set of depths
below \(s_{\max}\) can be used simultaneously on a further dense
restriction, with current individual floors, unchanged earlier
full-domain witnesses, and the original true assignment priors.
\end{proposition}

\begin{proof}
At \(u=1\), \Cref{a06:horizontal-coefficient-count} gives the
property by binned affine coefficients, with no \(P\)-dependence in
the predictors.  We establish a strict improvement for every
\(2k<u\le1\), and then construct the endpoint objects.

Fix \(2k<u\le1\), and suppose \(\mathcal G(u)\) holds.
For a target depth \(s_1<s_{\max}\), choose
\(s_0=s_1/(1+\theta)\), where the fixed positive \(\theta\)
will be specified below. Write
\[
 a_0=a_{s_0},\qquad D_0=a_0^u,\qquad
 W_G=\pi(d_0)D_0^2.
\]
Prepare the nearby true depths, and then take the full-time
predictors \(H_0\) supplied by \(\mathcal G(u)\) at \(s_0\).
\Cref{a07:full-time-comparison} compares these predictors with
the true affine entries. We will improve the fitted width
from \(a_0\) to \(a_{s_1}\) on the same parameter box.

\par\smallskip\noindent\textbf{The conditional prior and scalar support.}\quad  Choose one such paired \(L_{s_0},Q_{s_0}\) per \(H_0,I\) losing only subpower in aggregate success. Thus \({\cal L}_{s_0}\) is time-block determined within \(H_0\). Use typical \(H_0\) as charts with dense conditional success and conditional prior mass denominator
\[
 \nu^\flat(H_0)\sim_{\log,N} W_G a_0^d .
\]
Indeed restore floors or discard now-light \(H_0\) by the threshold
sums; the joint cap near its graph gives the matching incidence
upper bound. Prior-times-time domination and deletion of poor
success fractions, charged to the disjoint full-time assignments,
give the stated prior denominator. Fix one such chart and retain
the trajectory probability
\[
 \nu_H=\nu^\flat(\,\cdot\mid H_0)
\]
throughout this bootstrap step. Let \(\lambda_H\) be the current
incidence measure in this chart divided by \(\nu^\flat(H_0)\).
It is a submeasure with dense mass, not a replacement trajectory
prior. Write
\[
 d\lambda_H=f_H\,d\nu_H\,dt,\qquad
 0\le f_H\le K_{\rm pre},\qquad
 \delta_H=\lambda_H(\mathrm{all})\ge1/K,
\]
where \(K_{\rm pre}\) is subpower. In this chart put
\[
 y=(P-P_c)/D_0,\qquad \bar x=(x-H_0(t,P))/a_0 .
\]
Thus \(y\in\mathbb R^2\) is static, \(\bar x\) affine on lines with slope and range \(\le K\) by the full-time prediction. On active time blocks put \(F(t)=\nabla_y({\cal L}_{s_0}-H_0)/a_0\), a bounded-by-\(K\) time-only row by the comparison (it can be set to zero on unused blocks). Put \(S(t)=\bar x(t)-F(t)y\).

The immediate target is a short list of residual values and slopes
at one time. With \(w=N^{-\theta s_0}\), where \(\theta\) is
chosen below, we seek a bounded time-only row \(A(t)\) and a
chart-dependent time \(t_*\) such that on dense labels only
\(Kw^{-d}\) bin pairs
\[
 \bigl(S_l(t_*)_w,\ (\bar x_l'-A(t_*)y_l)_w\bigr)
\]
are needed. Since \(y_l\) is static and \(\bar x_l\) is
affine in time, each pair determines a full-time prediction affine
in \(y\), with affine time coefficients and error \(Kw\).
Restoring \(H_0\) gives
thickness \(Ka_0w=Ka_{s_1}\) on the unchanged parameter box:
\[
 D_0=a_0^u=a_{s_1}^{u/(1+\theta)}.
\]
This is how the step improves \(\mathcal G(u)\) to
\(\mathcal G(u/(1+\theta))\). The following preparations verify
the hypotheses of \Cref{prop:scalar-normal}, whose first-jet
conclusion supplies this value-and-slope list.

Take the accuracy exponent \(E=(u-2k)s_0/10,\ \zeta=N^{-E}\); in particular
\[
 D_0/h_{s_0+E}^2\ll\zeta.
\]
We first obtain integrated mixed caps for \(\lambda_H\), then
prepare time labels to obtain instantaneous caps under \(\nu_H\).
For widths \(r,w\in[\zeta,1]\), integrate the joint \(d_0\)-cap
at thickness \(a_0w\) about the \(H_0\) offset. On \(P\)'s
moving spatial domain the result, before division by the chart
denominator, is at most
\(K\pi(d_0)(D_0r)^2(a_0w)^d\) per unit time.
Subdivide \(P\) as finely as needed to evaluate that offset. The
absolute time and parameter meshes may have depths that are large
fixed multiples of \(s_{\max}\): the coefficient amplification
is at most \(K/(D_0a_0)\), and the reconstruction buffer controls
these errors. Still finer true meshes may be used for the
integration. Consequently the normalized bound over a tested time
bin \(I\) is \(K|I|r^2w^d\).

Set \(E_l^0=\{t:f_H(l,t)\ge\delta_H/2\}\), and let
\(d\rho_H^0=\mathbf1_{E_l^0}(t)d\nu_H(l)dt\). The omitted
\(\lambda_H\)-mass is at most \(\delta_H/2\), and on the
retained event
\[
 (\delta_H/2)\,d\rho_H^0
 \le d(\lambda_H|_{E^0})\le K_{\rm pre}\,d\rho_H^0.
\]
In particular \(\rho_H^0\) has mass at least
\(\delta_H/(2K_{\rm pre})\).

Choose one tiny time partition, fine enough for the current finite
width grid that \(K|I|\ll\zeta\). For each indexed line \(l\)
and time bin \(I\), retain \(E_l^0\cap I\) only if
\(|E_l^0\cap I|\ge |I|/K_{\rm dur}\). Let \(E_l\) be the
union of these retained sets and
\(d\rho_H=\mathbf1_{E_l}(t)d\nu_H(l)dt\). The deletion costs
at most \(1/K_{\rm dur}\) in indicator mass and
\(K_{\rm pre}/K_{\rm dur}\) in weighted mass. Choose this
subpower threshold small compared with the two dense masses.

For fixed \(t\in I\), let \(A_t\) be the trajectories currently
labeled at \(t\), with \(y\) in a specified \(r\)-square and
\(\bar x(t)\) in a specified \(w\)-interval. Bounded speed
puts their entire retained portion of \(I\) in the same enlarged
constraints. The duration floor and integrated cap therefore give
\[
 \frac{|I|}{K_{\rm dur}}\nu_H(A_t)
 \le \rho_H^0(I\times\text{enlarged constraints})
 \le \frac{2}{\delta_H}
       \lambda_H(I\times\text{enlarged constraints})
 \le K|I|r^2w^d.
\]
After division by \(|I|\), and absorption of subpower factors,
\[
 \nu_H\{l:t\in E_l,\ y_r,\bar x_l(t)_w
                    \text{ in specified bins}\}\le Kr^2w^d.
\]
These are line--time-bin deletions. All subsequent selections
restrict \(E_l\), keeping \(\nu_H\) fixed, so this instantaneous
upper bound persists. Null exceptional times and trajectories may
be omitted. The displayed density comparison also returns any
retained indicator mass to the weighted incidence law at subpower
cost.

We also have at each active \(t\), writing \(w=N^{-s'}\), \(0\le s'\le E\),
\[
 N_w\{S(t)\}\le K w^{-d}.
\]
This support count uses the earlier full-domain pair witnesses in
the horizontal law, frozen before conditioning to \(H_0\). Current
\(\rho_H\)-labels locate the contributing boxes; the witness
floors and their common-box ceiling are both measured in that
earlier law.
Use the true \(Q_s,\ s=s_0+s'\), paired with the shared \(Q_{s_0},L_{s_0}\) on this time block. Their domains meeting the occurrences at \(t\) are all in a common enlarged horizontal box per child time interval (as in the full-time count): contact error along the \(G\) reference motion costs \(K D_0\), now much smaller than \(h_s^2\) up to \(K\). Sum their earlier full-domain pair witnesses of individual mass \(\ge K^{-1}h_s^4 a_s^d\) in this common box, satisfying the same \(L_{s_0}\)-test. The pure cap near that fit gives a total ceiling \(K h_s^4 a_0^d\), hence \(\le K(a_0/a_s)^d\) possibilities, counting the affine subentries as well. On a pairing, \(S\) differs within \(K w\) from a common time-only term plus \(({\cal L}_s-{\cal L}_{s_0})/a_0\) by affine parameter dependence and the definition of \(F\). At fixed \(t\) this last function has parameter slopes bounded by \(K/h_s^2\): use \(\|L_s-L_{s_0}\|\le K a_0\) on the whole centered child box from the pair witnesses. Thus its variation across the \(D_0\)-box costs \(\le K w\). Evaluation errors using \Cref{a06:horizontal-equation} fit this estimate as well. This proves the support count.

At each tested width \(w\), discard scalar bins of instantaneous
\(\nu_H\)-labeled mass less than \(w^d/K\), with \(K\)
sufficiently enlarged. The support count, integrated in time, pays
for this deletion. To bound mass in an \(r\)-interval of \(S\),
\(w\le r\le1\), apply the mixed cap with \emph{scalar width
\(r\)} and parameter squares \(y_r\). On each such square the
bounded row \(F(t)\) turns the \(S\)-constraint into at most
\(K\) enlarged \(\bar x\)-bins of width \(r\). Summing the
\(Kr^{2+d}\) cap over at most \(Kr^{-2}\) parameter squares
gives \(Kr^d\). Division by the retained scalar-bin floor gives
the local count \(K(r/w)^d\).

Prepare the true-chart and pair witnesses on fixed finite exponent
grids before applying the baseline property to that dense input.
After obtaining \(H_0\), the preceding counts hold on these grids,
and the scalar-bin deletion costs a sum of arbitrarily small
reciprocal-subpower losses. Let the grids densify slowly enough
that their cardinalities and total losses remain subpower. Rounding
and interpolation then give the local counts and caps on the whole
width range. The next two preparations further restrict these
labels, first on finite grids and then on a slower diagonal.

\par\smallskip\noindent\textbf{Regularizing the scalar-normal input.}\quad  First we require the Lipschitz comparison for \(F\) on active times, with \(\zeta\)-slack. At a tested dyadic width \(r\), sample uniform time pairs from common \(r\)-intervals (averaged with length weights). The mean prior line weight jointly labeled there is at least an inverse-subpower by Jensen on the labeled time fractions in these intervals. But for \(|t-t'|\le r\) active with \(M=|F(t)-F(t')|\gg r\), this shared-line weight is \(\le K(r/M)^{1-d}\).

To prove this bound, given a scalar source \(S(t)_r\)-bin, target \(S(t')\)'s for common lines lie within \(K(M+r)\) of it since \(\bar x\) has bounded speed up to \(K\). There are \(\le K(M/r)^d\) target bin choices by local scalar counts. Thus the projection of static \(y\) onto the difference axis lies in that many intervals of length \(K r/M\), using the signal equation. Each projection interval costs at most \(K/(M r)\) squares of side \(r\) for \(y\) (\(M\le K\)). With the given scalar source bin each such square costs \(\le K r^{2+d}\) prior weight at \(t\). Sum over \(\le K r^{-d}\) source bins to obtain the estimate.

Consequently the pairs with \(M/r\) larger than a sufficiently large subpower contribute negligibly to the shared-label lower bound. By averaging choose one anchor time per interval; retaining labels at times with \(|F-F_{\rm anchor}|\le K r\) leaves dense total mass (one can count even just shared lines for this lower bound). Retain also a separated residue class of intervals with dense total mass, so that active times within distance \(r\) compare in the same interval. Iterate at all tested scales, upper bounds persisting under restriction, to obtain the comparison \(K(|t-t'|+\zeta)\) using slow grids.

Next prepare the slope comparison on spacetime \((t,y,\bar x)\)
and floors for coefficient bins under the fixed prior \(\nu_H\).
At width \(r\), a typical active \(y_r\)-group has prior and
label mass of order \(r^2\), up to \(K\), by the cap, summable
thresholds and deletion of poor success fractions. At fixed time
it has at most \(Kr^{-d}\) bins of \(\bar x_r\), by the
\(S\)-count and boundedness of \(F\).

For each line and \(r\)-time interval \(I_r\), omit the pair
if its current labeled duration is less than \(r/K\), with
\(K\) sufficiently large. Every now-occurring spacetime bin
has witnesses within its \(Kr\)-enlargement for duration at
least \(r/K\). Integrating the fixed-time support count over
these witnesses gives at most \(Kr^{-1-d}\) spacetime bins per
\(y_r\)-group. Also retain lines with total labeled duration
at least \(1/K\), and denote this set by \(\mathcal L_r\).
For clustering in a typical group, use the auxiliary finite prior
obtained by restricting \(\nu_H\) to that group and
\(\mathcal L_r\), and dividing by \(r^2\). Its mass is at most
\(K\). Testing a coefficient ball on those long labeled times
and using the integrated spatial caps bounds this auxiliary
prior's coefficient-ball masses by \(Kw^d\), for \(w\ge r\).
The reference probability \(\nu_H\) itself remains fixed.

\Cref{lem:scalar-clustering} now permits retaining joint bins
\(B_r=(y_r,[\bar x\text{ coefficients}]_r)\), each of current
\(\rho_H\)-mass at least \(r^{2+d}/K\). Indeed apply clustering
anew to any dense putative residual of lighter bins in typical
groups, with the same support and long-duration preparations. The
cluster's auxiliary prior floor, multiplied by the long-duration
threshold, gives a heavy ordinary active bin, a contradiction.
At reciprocal-subpower widths all parameter counts are already
subpower.

Next prune the distinct entries \((B_r,I_r)\), where \(I_r\)
is an \(r\)-time interval, below the current mass floor
\(|I_r|r^{2+d}/K\). The heavy-bin count makes these thresholds
summable. If such an entry meets a given spacetime \(r\)-cube,
its entire earlier subentry has \((y,\bar x)\) within \(Kr\)
of that cube: coefficients differ by \(O(r)\) inside \(B_r\),
and motion over \(I_r\) is bounded by \(Kr\). Summing its full
subentry floor under the integrated spatial cap gives at most
\(K\) possible coefficient bins, hence slope bins, for the cube.

In each cube select one slope bin with dense aggregate success,
then retain a separated residue class of cubes so that selected
points within distance \(r\) lie in the same cube. For example,
color all coordinate indices modulo a fixed integer greater than
two and keep a dense class. Iterating over the tested scales gives
the residual velocity comparison with \(\zeta\)-slack, since
\(y'=0\).

For each participating full parameter bin \(B_r\), the active
floor just obtained implies \(\nu_H(B_r)\ge r^{2+d}/K\),
since time has length one. This fixed-prior floor survives later
label restrictions. The absolute instantaneous cap therefore gives
\[
 \frac{\nu_H\{l\in B_r:t\in E_l,
                   (y_l,\bar x_l(t))\text{ in a }q\text{-cube}\}}
      {\nu_H(B_r)}
 \le K\frac{q^{2+d}}{r^{2+d}},\qquad \zeta\le q<r.
\]
Only static \(y\) and the affine \(\bar x\)-coefficients are
binned; hidden trajectory indices remain in \(\nu_H\). Nested
dyadic grids allow interpolation using a participating finer bin's
prior floor inside its containing bin. The finite iterations and
slow diagonal retain dense mass, with the finest-width comparisons
also controlling smaller distances. These are all the hypotheses
of \Cref{prop:scalar-normal}, under the unchanged prior \(\nu_H\)
and the current labels \(E_l\).

\par\smallskip\noindent\textbf{One bootstrap step.}\quad  Apply the first-jet conclusion of \Cref{prop:scalar-normal} to the normalized signals \(\bar x,y\). On dense labels in the chart we have slope row \(A(t)\) bounded by \(K\) and
\[
 \mathsf H((\bar x'-A(t)y)_w\mid t,S(t)_w)=o(\log N)
\]
for \(w=N^{-\theta s_0}\), taking \(\theta=(1-d)(u-2k)/10000\); here sample trajectory and time independently before imposing the retained labels. The finest affine time-fit exponent can be \(E/4\), with dyadic
rounding.  Indeed, for \(H_*=N^{-E/4}\), the floor satisfies
\(\zeta=N^{-E}\ll H_*^3\), and
\[
 \theta s_0=\frac{(1-d)E}{1000}
       <\frac{(1-d)E}{256}.
\]
These are the floor and reading conditions of the first-jet
conclusion.  Also
\(D_0/h_{s_0+E}^2=N^{-(10-2k)E}\ll\zeta\), so the preceding
support comparison has fixed positive-power room. All dense successes and subpower bounds can be taken uniform across the typical charts with the stated common inputs (or test against sequences of such charts). The negligible log entropy permits retaining, given time and scalar bin, subpower lists of the indicated residual bin (discard very small conditional probabilities). Thus at fixed active time \(\le K w^{-d}\) total bin pairs suffice, by the scalar support bound. Discard also lines with too short total labeled duration, then pick by averaging a single time \(t_*\) within the chart with dense line weight thus covered. These trajectories account for dense labeled incidence mass using their inverse-subpower labeled durations, also restoring earlier weights at subpower cost.

Assign these lines disjointly by the bin pair at \(t_*\), fixed within their earlier full-time \(H_0\) label. For bin centers \(b,b'\), predict throughout by
\[
 \bar x(t)= b+F(t_*)y+(t-t_*)\big(b'+A(t_*)y\big)+O(Kw)
\]
using exact affine evolution. Restoring \(H_0\) gives the prescribed predictor type bounded by \(K\) on the unchanged full-time scaled parameter box, at thickness \(a_0 w=a_{s_1}\). Typical heavy entries among the new labels have the required active mass lower \(\pi(d_0)D_0^2 a_{s_1}^d/K\): per earlier chart the prior denominator compares with \(W_G a_0^d\), and there are only \(K w^{-d}\) new labels, so drop light ones by summing. This proves \({\cal G}(u/(1+\theta))\) at the arbitrary target \(s_1<s_{\max}\). All parameter and time accuracies needing reconstruction here used only fixed multiples of \(s_{\max}\) in absolute depth, independently of how small \(u-2k>0\) is; each fixed application has positive fixed jet exponents before any endpoint diagonal.

\par\smallskip\noindent\textbf{A fixed reconstruction buffer.}\quad
The accuracy requirements in the preceding step are uniform in the
distance \(u-2k\) in the following precise sense.  The normalized
parameter and scalar widths are at least \(\zeta=N^{-E}\), with
\(E\le s_{\max}/10\).  The first-jet argument uses the fixed-order
fitting scales \(H^{1/4}\), \(H^{1/8}\),
\(H^{(1-d)/32}\), and interpolation widths above this floor.
As \(E\) decreases, these meshes become coarser.  The largest
coefficient amplification when returning the baseline normalized
predictor to physical parameter coordinates is
\[
 \frac{K}{D_0a_0}\le K N^{2s_{\max}},
 \qquad 2k<u\le1.
\]
All further differentiated tests have fixed degree and order.  Thus
one absolute exponent \(C_0s_{\max}\), for a fixed sufficiently
large \(C_0\), supplies the internal time and parameter meshes for
every finite bootstrap step.  Choose \(s_{\max}\) at the outset so
that \(C_0s_{\max}<e_0/2\), enlarging \(C_0\) first to include
the fixed reconstruction and evaluation amplifications.  The
\(N^{-e_0}\) reconstruction error is then below the meshes and
their required scalar accuracies by a fixed power.  The number of
tested widths, coloring losses, and finite-iteration constants may
depend on the step; they do not require an exponent proportional
to \(1/E\).  Finer true-base meshes used only to integrate an
already established cap do not invoke reconstruction at those
meshes.

\par\smallskip\noindent\textbf{The endpoint sequence and its actual labels.}\quad
Put \(\alpha=(1-d)/10000\).  The recurrence is
\[
 u_{j+1}=\frac{u_j}{1+\alpha(u_j-2k)},\qquad u_0=1.
\]
It satisfies
\begin{align*}
 u_{j+1}-2k
 &=\frac{(u_j-2k)(1-2k\alpha)}{1+\alpha(u_j-2k)}>0,\\
 u_j-u_{j+1}
 &=\frac{\alpha u_j(u_j-2k)}{1+\alpha(u_j-2k)}>0.
\end{align*}
The limit must therefore be \(2k\).  Fix a desired target depth
\(s<s_{\max}\) and a dense input sequence.  For each finite \(j\),
perform its finite sequence of bootstrap steps on a further
subsequence.  Absorb every loss, grid count and norm bound in
\(K_j(N)\).  Take \(N\) sufficiently late on that subsequence that
\(\log K_j(N)/\log N<1/j\), that the retained relative mass is
at least \(N^{-1/j}\), and that every required finite accuracy
buffer holds.  A diagonal with \(j=j(N)\to\infty\) chosen only
after these thresholds has subpower total losses.  This construction
is made for the given input sequence; it asserts no uniform diagonal
over all possible dense inputs.

Let \(D_j=a_s^{u_j}\), \(D_*=a_s^{2k}\), and
\(R_j=D_*/D_j\).  Then
\(\log_N R_j=s(u_j-2k)\to0\).  Keep the old predictor labels
separate, even if several lie in the same new \(D_*\)-square,
and split each assignment by that new grid.  An old square meets
only boundedly many new squares.  Since the predictor is affine
in the two parameter coordinates, its norm on a meeting new square
over unit time is at most
\(C K_j(1+R_j)\).  Predictions on its assigned trajectories
are unchanged.

For each old label let
\(T_j=\pi(d_0)D_j^2a_s^d\).  Its earlier active floor gives
\(\sum T_j\le K_j\).  The required endpoint numerator for
each split label is \(R_j^2T_j\); hence the sum of these new
numerators is at most \(C K_jR_j^2\).  This is subpower.  Choose
the new floor constant larger than that sum divided by a small
inverse-subpower fraction of the current total mass.  Removing
light split labels keeps dense mass and yields precisely the
\(\pi(d_0)D_*^2a_s^d/K\) floor.  Both its norm bound and its
floor are therefore actual endpoint objects, even if \(R_j\)
itself tends to infinity.

For several prescribed depths, apply this construction successively
to the current dense input, preserving earlier assignments,
predictors and witnesses.  Each resulting family's floor
numerators have subpower sum.  Include also the full-time chart labels, with floor numerators
\(\nu^\flat(H)\) from their unchanged trajectory assignments.
These numerators sum to at most one per family by disjointness.
Apply \Cref{a07:simultaneous-floors} to all the families together.
The final law therefore has both the required predictor floors
and current chart mass at least \(\nu^\flat(H)/K\), despite
cascading deletions. Prior-times-time domination supplies the reverse
comparison. No further deletion is needed to obtain the conditional
prior denominators, and no earlier full-domain witness or true prior
is replaced by the final restricted law.  This proves \(\mathcal G(2k)\) with the stated
simultaneous-use assertion.
\end{proof}

\subsection{Returning to wide intervals away from the critical rate}

\begin{proposition}[Return to wide intervals]
\label{prop:offcritical-wide}
If \(2k\ne1\), every substantial residual of the horizontal
problem supplies the hypotheses of \Cref{prop:wide-interval} at
three fixed nearby depths.  Consequently it yields a
time-improvement candidate with a fixed positive depth gain, full
block fit, hidden derivative one, and the required original
true-descendant mass.  The assertion holds both for \(2k<1\)
and for \(2k>1\).
\end{proposition}

\begin{proof}
\par\smallskip\noindent\textbf{Parameter charts and their unchanged prior denominator.}\quad


Keep the horizontal-model coordinates, relative depths and true
assignment priors fixed. We will construct projected labels at three
nearby depths whose common transverse width is
\(D_0=h_{s_0}^2\). Put
\[
 D_0=N^{-2k s_0},\qquad \widetilde P=(P-P_c)/D_0
\]
using groups at \(P_{D_0}, d_0=[v]_\chi\), \(P_c=(p_c,q_c)\) the parameter-bin center (dyadic roundings as usual), common fine direction width \(\chi\le N^{-e_0}\). Take \(s_0>0\) sufficiently small compared with \(e_0\) for the fixed multiples below, also \(\tfrac32s_0<s_{\max}\) when using \({\cal G}(2k)\). Depth choices below can be fixed in advance including sufficiently fine \(\chi\) for the wide comparisons.

Prepare the earlier affine data (including \Cref{a06:affine-slope-bounds}) for the needed depths first. In particular for pairs used between \(s_0\) and a deeper \(s_j\) we can keep joint whole-domain witnesses of mass at least \(m_{Q_{s_j}}/K\) for the true labels and paired entries \(L_{s_j},L_{s_0}\). Use \Cref{a07:simultaneous-floors} to discard light such pairs before smaller slices/full-time selections, paying by the true child budgets and subentry counts; \(L_{s_j}-L_{s_0}\) then costs \(K a_{s_0}\) in norm throughout the child box. Whole-domain witnesses of this mass order for just the deeper labels, in the bounded \(x\)-problem here, suffice where no ancestor comparison is needed.

If \(2k>1\), use the groups \((d_0,P_{D_0})\) directly as charts with baseline predictor \(H_0=0\), thickness unit \(A_0=1\). If \(2k<1\), use \Cref{prop:bootstrap} and \Cref{a07:simultaneous-floors} to take simultaneous \(\mathcal G(2k)\) labels/predictors \(H_0,H_b\) at \(s_0,s_b=\tfrac32s_0\), with the same \(d_0\)'s. Compare each with \(L_{s_j},Q_{s_j}\) at its own depth by \Cref{a07:full-time-comparison} (in both cases write \({\cal L}_{s_j}(t,P)=L_{s_j}(t,p+t v_{d_0},q_0-tp)\) at either depth). Thus for \(2k<1\) the discrepancy of \({\cal L}_{s_j}\) to its paired full-time predictor at \(s_j=s_0,s_b\) costs \(K a_{s_j}\) on that time block times parameter domain of width \(N^{-2ks_j}\) up to \(K\). In this case take one paired \(L_{s_0},Q_{s_0}\) per \(H_0\) and time block using the list bound, and use \(H_0\) labels as charts with \(A_0=a_{s_0}\).

In either case keep typical such charts with active mass and prior mass denominator both of order up to \(K\)
\[
 W_G A_0^d,\qquad W_G=\pi(d_0)D_0^2 .
\]
Indeed the thresholds sum to subpower by the earlier label floors (or \(\sum W_G\le K\) over occurring groups in the trivial case). Incidence mass per chart costs at most this order by the joint \(d_0,x\)-cap near \(H_0\) at \(A_0\) using the moving parameter domain, as in the comparisons (\(\chi+N^{-e_0}\) errors negligible here). And one can discard poor success fractions relative to trajectory assignment weights in \(\nu^\flat\), by disjointness; we have prior-times-time domination in the upper direction as well. Condition the prior to each chart and divide active masses there by the prior denominator. If trimming other lines of the assignments is needed, e.g. to use bounded \(x\), one can do it at subpower cost since all our incidences already meet the trajectory bounds.

\par\smallskip\noindent\textbf{A scalar parameter and its fiber direction.}\quad
The scalar projection will use
\[
 R(t)=\widetilde P_2-2t\widetilde P_1,
 \qquad W(t)=(1,2t).
\]
At fixed time the fibers of \(R\) in the parameter plane have
direction \(W(t)\). We seek a correction after which the affine
parameter predictors vary by no more than their fitted width along
these fibers. They can then be evaluated at \(\widetilde P=(0,R)\)
and counted using only \((t,R)\).

Suppose first that \(2k<1\). The bounded-by-\(K\) time signal
on paired events in a chart
\[
 G_0(t)=(D_0/A_0)\nabla_P({\cal L}_{s_0}-H_0)\cdot W(t)
\]
agrees with \((D_0/A_0)\nabla_P(H_b-H_0)\cdot W(t)\) there to error
\[
 K\big(D_0/h_{s_b}+(D_0/A_0)a_{s_b}/h_{s_b}^2\big)\le N^{-\kappa}
\]
with, for example, \(\kappa=\tfrac{s_0}{4}\min\{k,1-2k\}>0\) after absorbing subpower losses. Indeed \(\nabla_P{\cal L}_j\cdot W=\partial_Z L_j+t\partial_Y L_j\). The first error compares \(L_{s_b}\) and \(L_{s_0}\) by the whole-child-box bound \(K A_0\): the slopes of the difference along the central horizontal plane cost \(K A_0/h_{s_b}\), transversely \(K A_0/h_{s_b}^2\), and the horizontal spatial displacement \((1,t)\) at the occurrence differs from \((1,t_c)\) of that plane by at most \(K h_{s_b}\). The second error compares the finer calligraphic fit with \(H_b\) using norm closeness throughout its paired parameter domain. The savings before subpowers are \(N^{-ks_0/2}\) and \(N^{-(1-2k)s_0/2}\) respectively.

We now choose one such quadratic time polynomial while retaining
mass from all paired trajectories. In a typical chart let
\(\lambda_{{\rm pair},H}\) be the incidence measure after the
pairings just used, divided by the chart's prior denominator, and
let \(\nu_H\) be that fixed conditional trajectory probability.
Thus \(\lambda_{{\rm pair},H}\le K\nu_H\otimes dt\). Denote
its time marginal by \((\lambda_{{\rm pair},H})_t\), and write
\[
 m_H(t)=\frac{d(\lambda_{{\rm pair},H})_t}{dt},\qquad
 M_H=\int_0^1m_H(t)\,dt\ge1/K.
\]
Retain \(T_0=\{t:m_H(t)\ge M_H/2\}\). This discards at most
\(M_H/2\) of the paired incidence mass. Domination by
\(K\nu_H\otimes dt\) then supplies an indexed trajectory whose
paired time set \(T_b\subset T_0\) has
\(|T_b|\ge M_H/(2K)\). Its label \(H_b\) is fixed along the
whole trajectory.

For this fixed label,
\((D_0/A_0)\nabla_P(H_b-H_0)\cdot W(t)\) is a quadratic in
time. It approximates the shared signal \(G_0(t)\) to
\(KN^{-\kappa}\) on \(T_b\), so boundedness of \(G_0\) and
Remez give coefficient bounds \(K\). Restore every paired line
at those same times. The \(s_0\) entry was chosen per chart and
time block, so \(G_0\) is shared by all of them, and the restored
mass is
\[
 \lambda_{{\rm pair},H}\{t\in T_b\}
 =\int_{T_b}m_H(t)\,dt
 \ge (M_H/2)|T_b|\ge M_H^2/(4K).
\]
It is therefore dense. This argument uses the selected line to find
the time polynomial; the retained mass is that of all paired lines
on its time set.

Choose \(H_{\rm corr}(t,\widetilde P)\), affine in parameters
with affine time coefficients and bounded by \(K\) on the scaled
domain, whose derivative along \(W(t)\) is this polynomial.
For \(b_0+b_1t+b_2t^2\), take
\((b_0+b_1t)\widetilde P_1+(b_2t/2)\widetilde P_2\).
When \(2k>1\), set \(H_{\rm corr}=0\).

\par\smallskip\noindent\textbf{Choosing all three depths with positive margins.}\quad
Let \(C_d=10/(1-d)\), as in \Cref{prop:wide-interval}, and put
\(C=C_d+3\).  Reduce \(s_0\) so that \(100s_0<e_0\), all
preceding fixed reconstruction buffers hold, and
\(3s_0/2<s_{\max}\) in the subcritical case.  Set
\[
 B_*=
 \begin{cases}
 \min\{s_0/2,\kappa,ks_0/(1+k)\},&2k<1,\\
 \min\{s_0/2,(2k-1)s_0\},&2k>1.
 \end{cases}
\]
Choose
\[
 \delta=B_*/4,\quad s=s_0+\delta,\quad
 u=\min\{\delta/4,k\delta/(4C)\},\quad
 c_2=(1-d)u/8,
 \qquad s_i\in\{s,s+u,s+c_2\}.
\]
All these are fixed positive exponents.  Writing
\(\Delta_i=s_i-s_0\), we have
\(0<\Delta_i\le5B_*/16<B_*\), so \(s_i<3s_0/2\) and
\(D_0/h_{s_i}=N^{-k(s_0-\Delta_i)}\ll1\).  For \(2k<1\),
\begin{align*}
 \frac{D_0/h_{s_i}}{a_{s_i}/A_0}
 &=N^{-[ks_0-(k+1)\Delta_i]}\ll1,\\
 \frac{N^{-\kappa}}{a_{s_i}/A_0}
 &=N^{-(\kappa-\Delta_i)}\ll1.
\end{align*}
For \(2k>1\),
\[
 \frac{D_0}{a_{s_i}}
   =N^{-[(2k-1)s_0-\Delta_i]}\ll1.
\]
These inequalities supply all the directional error margins.  They
also give \(0<c_2<(1-d)u/4\), as required for the wide-interval
argument. On the paired events
\begin{equation}
 \left|\nabla_{\widetilde P}\big(({\cal L}_{s_i}-H_0)/A_0-H_{\rm corr}\big)\cdot W(t)\right|
          \le K a_{s_i}/A_0 .                                       \label{a07:directional-collapse}
\end{equation}
For \(2k>1\) use the horizontal bound \Cref{a06:affine-slope-bounds} times \(D_0\). For \(2k<1\) compare the true affine fits \(L_{s_i},L_{s_0}\) on the child box in the same horizontal direction as just used (cost \(K D_0/h_{s_i}\) here), then use the time-signal approximation. These derivatives only concern affine dependence on parameters.

\par\smallskip\noindent\textbf{The scalar projection: measure ceilings and support.}\quad
Apply \Cref{thm:scalar-mesh-projection} to the scalar parameter
\(R\) defined above and the signal
\[
 x^\circ=(x-H_0(t,P))/A_0-H_{\rm corr}(t,\widetilde P),\qquad
 w_i=a_{s_i}/A_0.
\]
Both signals are affine on trajectories with coefficients \(\le K\) in the conditional charts, by the full-time bounds. We check the scalar mesh hypotheses there at terminal width \(w_i\), with density prefactor one and angular gain, say using \((t,R)\) side \(\sigma_R\) of depth \(20s_0\). At fixed time the map from \(\widetilde P\) to \((R,R')\) has determinant of magnitude \(2\). Before division by the chart denominator, the mass with \((t,R)\) in a cell, \(x^\circ\) in a bin of width \(1\ge\varrho\ge w_i\), and optionally \(R'\) in an interval of width \(N^{-v_0}\), \(0<v_0\le s_0\), costs at most
\[
 K\pi(d_0)\,\sigma_R^2 D_0^2 (A_0\varrho)^d
\]
times \(N^{-v_0}\) if using that interval. Indeed reconstruct \(P\) from \((Z,Y)\) with error \(K(N^{-e_0}+\chi)\), using the determinant-one map \(P\mapsto (p+t v_{d_0},q_0-tp)\). Errors even after division by \(D_0\) are negligible compared with the cell precisions and \(a_{s_i}\). Integrate the true joint cap in much finer \((t,Z,Y)\) cells at fixed \(d_0\), with spatial area cost \(K D_0^2\sigma_R\) at each such time (times the further width factor if using \(R'\)). Test \(x\) to width \(K A_0\varrho\) with offset using reconstructed \(H_0+A_0 H_{\rm corr}\); coefficients in the indicated scaled parameters cost \(K\), so this uses only subpower lists per fine base cell. Thickening by sufficiently fine fixed-power true meshes changes the volume bound at most by \(K\). This proves the mass ceilings after division by \(\sim_{\log,N} W_G A_0^d\).

At each \((t,R)\)-cell the relevant \(Q_{s_i}\) domains lie in a common enlarged horizontal box per time block (the cell meets only \(K\) blocks). Specifically for the bin representative \(R_*\) use reference motion
\[
 Z_{\rm ref}=p_c+t v_{d_0},\qquad Y_{\rm ref}=q_c+D_0 R_*-t p_c .
\]
At contact their spatial discrepancies are \(D_0\widetilde P_1(1,t)\) up to \(K(D_0\sigma_R+N^{-e_0}+\chi)\) errors. The reference motion is exactly horizontal. Use its point/plane at block midpoint: \(D_0\ll h_{s_i}\) and the contact errors fit inside \(h_{s_i}^2\), with slope change \(K h_{s_i}\) for \((1,t)\). The aligned true boxes extend from these contacts with horizontal size \(K h_{s_i}\), transverse size \(K h_{s_i}^2\) relative to planes differing by at most \(K h_{s_i}\) in slope there. Thus the common enlarged box of base volume \(K h_{s_i}^4\) suffices for their whole domains.

The common box bounds the number of contributing true labels,
including their affine entries, by \(Kw_i^{-d}\). Indeed for \(2k>1\) sum the disjoint whole-domain witness floors \(m_{Q_{s_i}}/K\) by the base cap alone (bounded \(x\) and unit cap). For \(2k<1\) these earlier paired witnesses lie near the shared \(L_{s_0}\) of the chart/time-block selection at \(s_0\), so their summed weight costs \(\le K A_0^d h_{s_i}^4\) by the pure \(x\)-band cap. These counts use witnesses and caps in the full horizontal problem before the parameter-chart conditionings.

Each affine entry yields only \(K\) bins for \(x^\circ\) at \(w_i\) on the cell. It predicts using \(({\cal L}_{s_i}-H_0)/A_0-H_{\rm corr}\) with error \(K w_i\); replacing the true spatial variables as in this formula is allowed by \Cref{a06:horizontal-equation,a06:affine-slope-bounds}. Moreover one can evaluate this expression at \(\widetilde P=(0,R(t))\) instead by \Cref{a07:directional-collapse} and bounded \(\widetilde P_1\). To control its variation within the cell, use
\Cref{a06:affine-slope-bounds} and an active position: the ordinary
slopes of \(L_{s_i}\) are at most
\(K(1+a_{s_i}/h_{s_i}^2)\le KN^{(2k-1)_+s_i}\).
The functions \(H_0,H_{\rm corr}\) have norm \(K\) on the
scaled parameter domain and its bounded enlargements.  After division
by \(A_0\), the evaluated entry's time and \(R\) derivatives
are bounded by \(KN^{5s_0/2}\), using \(s_i<3s_0/2\) and
\(0<k\le1\).  Its variation across a \(\sigma_R\)-cell is
therefore at most \(KN^{-35s_0/2}\ll w_i\), because
\(w_i\ge N^{-3s_0/2}\).  Furthermore,
\[
 D_0\sigma_R\ll h_{s_i}^2,\qquad
 \frac{N^{-e_0}+\chi}{D_0}\ll\sigma_R
\]
when \(100s_0<e_0\).  Enlarging the fixed preliminary buffer also
covers division by \(A_0\) and all scalar evaluation errors.
Thus the projection uses reconstruction only at accuracies
strictly coarser than its known error. Each of at most
\(K\sigma_R^{-2}\) cells has at most \(Kw_i^{-d}\) terminal
scalar bins, as required.

\Cref{thm:scalar-mesh-projection}, with its angular gain, gives disjoint full-time trajectory labels of active mass \(\ge K^{-1}w_i^d\) in conditional units, predicting \(x^\circ\) throughout to \(K w_i\) by quadratics \(F_i(t,R)\) with coefficients \(\le K\). Indeed their interval width in that theorem and its reciprocal cost \(K\); the interval center and slope then cost \(K\) on occupied labels, so its moving-interval norm bound suffices. Retain simultaneous data at the three depths by applying successively to dense restrictions with the same ceilings/support bounds. Keep the ensemble of typical charts using the uniform subpower conclusions for the conditional charts, not just one frozen \(d_0\); thus there is still dense total mass in the horizontal problem.

\par\smallskip\noindent\textbf{Returning the predictors and their masses to true coordinates.}\quad  Split \(p=Z(0)\) to width \(\chi\) to return to \(\alpha\) for wide intervals, with representative \(p_{\alpha}\). Use \(P^\#=(p_\alpha,Y+t p_\alpha)\) instead of \(P\), and corresponding \(\widetilde P^\#,R^\#\), in
\[
 H_0 + A_0 H_{\rm corr}+A_0 F_i(t,R).
\]
The resulting predictor for \(x\) is quadratic in \(t,Y\) since \(P^\#,\widetilde P^\#,R^\#\) now depend affinely on those variables (the first scaled-parameter component is constant). It has norm \(K\) on the relevant moving interval of width \(K D_0\) about \(q_c-t p_\alpha\). Throughout the assigned lines we can use error \(K a_{s_i}\), since \Cref{a06:horizontal-equation} holds throughout participating lines and \(\|\widetilde P^\#-\widetilde P\|\le K(N^{-e_0}+\chi)/D_0\); trim lines to the trajectory bounds on counted incidences if needed. Thus the packet width scale for \(Y\) can be \(D_0\).

Their typical sliced weights are \(\ge K^{-1}a_{s_i}^d w_\alpha D_0,\ w_\alpha=\pi(d_0)\chi\). Before this split the floor per projected label (on obtaining that projection) is \(\pi(d_0)D_0^2 A_0^d w_i^d/K\); assignments at a depth in and between charts are disjoint. Remove light slices by summing thresholds over all those labels and the \(\le K D_0/\chi\) options each. Explicitly, an unsliced floor numerator is
\[
 T_i=\pi(d_0)D_0^2 A_0^d w_i^d
     =\pi(d_0)D_0^2a_{s_i}^d.
\]
The numerator required for one slice is
\(t_i=\pi(d_0)\chi D_0a_{s_i}^d\).  Since there are at most
\(KD_0/\chi\) slices, their total numerator is at most \(KT_i\).
The old floors and disjoint assignments give a subpower sum of
the \(T_i\)'s.  A larger common floor constant therefore pays the
new light-slice deletion.

Before the final common pruning, prepare the conditioning states
needed for \Cref{prop:wide-interval}. Each state consists of a true
\(Q_s\) and sufficiently fine true-base and \(x\)-bins; no
projected predictor label is appended. By \Cref{lem:horizontal-affine},
there are \(K\) choices per sufficiently deep end history \(e\).
Its earlier reference weight \(R_e\) bounds the velocity numerator
by \(K\pi(b)R_e\), and the references sum inside \(Q_s\)
to at most \(Km_{Q_s}\). Thus the numerators for all appended
history--state entries have total at most
\(K\sum_{Q_s}m_{Q_s}\le K\).

Apply \Cref{a07:simultaneous-floors} to these entries and all
three sliced projected layers together. This restores the projected
floors while giving each retained history--state entry mass at least
\(R_e/K\). Division of its persistent velocity numerator, then
summation over histories, gives conditional domination by \(K\pi\)
on the final law. All earlier whole-domain witnesses remain available.

Width \(D_0\) fits the wide condition since \(b_s/D_0\sim_{\log,N}N^{-2k(s-s_0)}\), \(D_0\ll h_s=g_s\); the chosen parameters satisfy
\[
 (C_d+3)u\le k\delta/4<\min\{2k(s-s_0),ks\}.
\]
Consequently both wide inequalities hold with a fixed power of
slack: \(b_s/D_0\) decays with exponent \(2k(s-s_0)\), and
\(h_s\beta_{\max}/\beta_{\min}\) decays with exponent \(ks\),
since all three widths \(\beta_i\) equal \(D_0\) up to
subpowers.

We finish by checking the wide-interval input and returning its
output to the initial tangent parent.

First, use \(H_s\) as the block length called \(q_s\) in
\Cref{prop:wide-interval}. The velocity \(v\) is unchanged by
the horizontal normalizations, so the same palette and interval
caps apply. The current dense law has the pure and joint caps in
\Cref{a06:residual-caps}. The common pruning just performed gives conditional velocity
domination at the fine true-base and \(x\)-states, also for offsets
of fixed-power variation in the true base. Selecting a heavy
descendant does not change these conditional laws, since \(Q_s\)
is already part of the state.

Second, each true label has only \(K\) affine entries \(L_s\)
predicting \(x\) to \(Ka_s\). Their norms in the centered
\(Q_s\) boxes and the needed enlargements are at most \(K\):
use an active value of \(x\), the horizontal slope bound in
\Cref{a06:affine-slope-bounds} multiplied by \(h_s\), and the
transverse slope bound multiplied by the transverse width
\(h_s^2\). The aligned planes
therefore give the required true tilted domains. The wide argument
applies with these lists and the \(\Phi\)-labels on true base
positions, including its final fine projection mesh.

Third, \Cref{prop:wide-interval} gives a quadratic prediction on a
true relative descendant \(Q_s\), with horizontal-problem mass
at least \(H_s\nu^\flat(Q_s)N^{-dc_2}/K\). Let
\(q_{\rm desc}\) and \(\nu_{\rm desc}\) be that descendant's
block length and full assignment weight in the initial tangent
parent. The time change and conditional full-assignment prior in
\Cref{lem:normalization} give
\[
 q_{\rm desc}=q_rH_s,\qquad
 \nu_{\rm desc}=\nu_{Q_r}\nu^\flat(Q_s).
\]
The additional restrictions and bounded-trajectory normalizations
already used affect comparisons of retained incidence mass by only
\(K\). Restoring the factor \(q_r\nu_{Q_r}\) therefore gives
counted mass at least
\[
 \frac{q_r\nu_{Q_r}\,H_s\nu^\flat(Q_s)N^{-dc_2}}{K}
 =\frac{q_{\rm desc}\nu_{\rm desc}N^{-dc_2}}{K}.
\]
Finally restore \(w=L+a_rx\). The resulting test has error
\(Ka_ra_{s+c_2}\) throughout the full assigned block.
Its hidden derivative is one, and its base terms are affine or
quadratic as required.

This is a candidate of the precise kind used in
\Cref{prop:candidate-upgrade}.  Its mass is charged to the original
residual and its fit is on the full assigned block.  The argument
works anew on each substantial residual, so the candidate-upgrade
criterion applies.
\end{proof}

Together \Cref{prop:critical-rate,prop:offcritical-wide} finish the
tangent contradiction for \(0<\ell<\infty\).
